\subsection{部分算子的例子}

在本节中， K = R 或 C 。

例. --- 1) 设 X 是局部紧拓扑空间， \(\mu\) 是 X 上的一个正测度。我们固定元素 \(p_1\) 与 \(p_2\) ，它们属于 \([1, +\infty[\) 。

设 \(g\) 是 \(\mu\)-可测的函数，定义在 \(X\) 上，取值于 \(K\) 中。用 \(D\) 表示 \(\mathcal{L}_{K}^{p_1}(X, \mu)\) 中满足如下条件的函数 \(f\) （属于 \(\mathcal{L}_{K}^{p_1}(X, \mu)\) ）所成的子空间： \(gf \in \mathcal{L}_{K}^{p_2}(X, \mu)\) 。由 \(D\) 到 \(\mathcal{L}_{K}^{p_2}(X, \mu)\) 中的、由 \(f \mapsto gf\) 所定义的线性映射确定了 \(\mathcal{L}_{K}^{p_1}(X, \mu)\) 到 \(\mathcal{L}_{K}^{p_2}(X, \mu)\) 中的一个部分算子，我们把它记作 \(m_g\) 。

\(\mu\)-可忽略函数在 \(\mathcal{L}_{\mathrm{K}}^{p_1}(\mathrm{X},\mu)\) 中所成的向量子空间包含在 D 中；并且， \(m_g\) 把 \(\mu\)-可忽略函数映为仍为 \(\mu\)-可忽略的函数。我们将用 \(\tilde{m}_g\) 表示从 \(\mathrm{L}_{\mathrm{K}}^{p_1}(\mathrm{X},\mu)\) 到 \(\mathrm{L}_{\mathrm{K}}^{p_2}(\mathrm{X},\mu)\) 中的、由 \(m_g\) 通过过渡到商诱导出的部分算子。它称为以 \(g\) 为乘子的乘法算子（从 \(\mathrm{L}_{\mathrm{K}}^{p_1}(\mathrm{X},\mu)\) 到 \(\mathrm{L}_{\mathrm{K}}^{p_2}(\mathrm{X},\mu)\) 中）。凡是 \(g_1\) 与 \(g_2\) 局部 \(\mu\)-几乎处处相等的函数，都定义同一个乘法算子。

命题 5. — 乘法算子 \(\widetilde{m}_{g}\) （从 \(\mathrm{L}_{\mathrm{K}}^{p_{1}}(\mathrm{X},\mu)\) 到 \(\mathrm{L}_{\mathrm{K}}^{p_{2}}(\mathrm{X},\mu)\) ）是稠定的闭算子。

我们先证明部分算子 \(\widetilde{m}_g\) 是闭的。设 \((f_n, h_n)_{n \in \mathbf{N}}\) 是 \(\mathcal{L}_{\mathrm{K}}^{p_1}(\mathrm{X}, \mu) \times \mathcal{L}_{\mathrm{K}}^{p_2}(\mathrm{X}, \mu)\) 中的一个序列，使得 \((\widetilde{f}_n, \widetilde{h}_n)\) 的类所作成的序列（即 \(f_n\) 与 \(h_n\) 的类）属于 \(\widetilde{m}_g\) 的图像，且当 \(\mathrm{L}_{\mathrm{K}}^{p_1}(\mathrm{X}, \mu) \times \mathrm{L}_{\mathrm{K}}^{p_2}(\mathrm{X}, \mu)\) 趋于无穷时在 \(n\) 中收敛。设 \((f, h)\) 是 \(\mathcal{L}_{\mathrm{K}}^{p_1}(\mathrm{X}, \mu) \times \mathcal{L}_{\mathrm{K}}^{p_2}(\mathrm{X}, \mu)\) 中的一对元素，使得它们的类所成的对 \((\widetilde{f}, \widetilde{h})\) 是 \((\widetilde{f}_n, \widetilde{h}_n)\) 的极限。

存在序列 \((f_{n_{k}})_{k\in\mathbf{N}}\) （从序列 \((f_{n})_{n}\) 中抽取的）使得 \(f_{n_{k}}(x)\) 收敛到 \(f(x)\) ，对 \(\mu\)-几乎所有的 x 成立（INT, IV, p. 131, § 3, no. 4，定理 3）。由此推出， \(h_{n_{k}}(x)=g(x)f_{n_{k}}(x)\) \(\mu\)-几乎处处收敛于 \(g(x)f(x)\) 。此外，序列 \((h_{n_{k}})\) 在空间 \(\mathcal{L}_{\mathrm{K}}^{p_{2}}(\mathrm{X},\mu)\) 中收敛于 h 。因此，函数 h 与 gf \(\mu\)-几乎处处相等（loc. cit.）。于是 \(\widetilde{f}\) 属于 \(\widetilde{m}_{g}\) 的定义域，且 \(\widetilde{h}=\widetilde{m}_{g}(\widetilde{f})\) 。这就证明了 \(\widetilde{m}_{g}\) 是闭的。

我们来证明 \(\widetilde{m}_g\) 的定义域在 \(\mathrm{L}_{\mathrm{K}}^{p_1}(\mathrm{X},\mu)\) 中稠密。只须验证：函数 \(f\in \mathcal{K}(\mathrm{X};\mathrm{K})\) 的类属于 \(\widetilde{m}_g\) 的定义域在 \(\mathrm{L}_{\mathrm{K}}^{p_1}(\mathrm{X},\mu)\) 中的闭包。设 \(f\in \mathcal{K}(\mathrm{X};\mathrm{K})\) ，并用 \(\widetilde{f}\) 表示它在 \(\mathrm{L}_{\mathrm{K}}^{p_1}(\mathrm{X},\mu)\) 中的类。对每个整数 \(n\in \mathbf{N}\) ，用 \(\varphi_{n}\) 表示由元素 \(x\in \mathrm{X}\) （满足 \(|g(x)|\leqslant n\) ）所成集合的特征函数，并令 \(f_{n} = f\varphi_{n}\) 。这时我们有 \(|gf_n|\leqslant n|f|\) ，而它属于 \(\mathcal{L}_{\mathrm{K}}^{p_2}(\mathrm{X},\mu)\) ，因此 \(f_{n}\) 属于 \(\widetilde{m}_g\) 的定义域。对 X 的每个元素 \(x\) ，序列 \((f_n(x))_{n\in \mathbf{N}}\) 收敛于 \(f(x)\) （当 \(n\to +\infty\) 时）；此外，我们有 \(|f_n|\leqslant |f|\) ，它属于 \(\mathcal{L}_{\mathrm{K}}^{p_1}(\mathrm{X},\mu)\) 。由勒贝格定理（INT, IV, p. 137, § 3, n°7，定理 6）， \(f_{n}\) 的类的序列在 \(\widetilde{f}\) 中收敛于 \(\mathrm{L}_{\mathrm{K}}^{p_1}(\mathrm{X},\mu)\) 。因此， \(f\) 的类属于 \(\widetilde{m}_g\) 的定义域的闭包。

在下述命题中，我们假设 \(p_1 = p_2 = 2\) 。

命题 6. --- a) 设 \(g'\) 是 X 上的 \(\mu\)-可测函数，满足 \(|g|\leqslant |g'|\) 。我们有 \(\mathrm{dom}(\widetilde{m}_{g'})\subset\mathrm{dom}(\widetilde{m}_{g})\) ，且 \(\mathrm{dom}(\widetilde{m}_{g'})\) 是部分算子 \(\widetilde{m}_{g}\) 的一个核；

\begin{enumerate}
\def\labelenumi{\alph{enumi})}
\setcounter{enumi}{1}
\tightlist
\item
  设 F 是 \(\mathcal{L}_{\mathrm{K}}^{2}(\mathrm{X},\mu)\) 的一个子空间，它与 \(\mathcal{K}(\mathrm{X};\mathrm{K})\) 的交在 \(\mathcal{K}(\mathrm{X};\mathrm{K})\) 中稠密，且它在 \(\mathrm{L}_{\mathrm{K}}^{2}(\mathrm{X},\mu)\) 中的像 G 包含在 \(\mathrm{dom}(\widetilde{m}_g)\) 中。若 \(|g|^2\) 局部 \(\mu\)-可积，则 \(\mathcal{K}(\mathrm{X};\mathrm{K})\) 包含在 \(\widetilde{m}_g\) 的定义域中，且 G 是 \(m_g\) 的一个核。
\end{enumerate}

我们来证明 \(a\) )。若 \(f \in \mathcal{L}_{\mathrm{K}}^{2}(\mathrm{X}, \mu)\) 属于 \(m_{g'}\) 的定义域（从而 \(fg' \in \mathcal{L}_{\mathrm{K}}^{2}(\mathrm{X}, \mu)\) ），则假设蕴含 \(fg \in \mathcal{L}_{\mathrm{K}}^{2}(\mathrm{X}, \mu)\) ，由此即得结论。

我们来证明 \(\mathrm{dom}(\widetilde{m}_{g'})\) 是 \(\widetilde{m}_g\) 的核，也就是说， \(\widetilde{m}_{g'}\) 的定义域在希尔伯特空间 \(\mathrm{E}_{\widetilde{m}_g}\) 中稠密。设 \(h \in \mathcal{L}_{\mathrm{K}}^2(\mathrm{X}, \mu)\) ，其类 \(\widetilde{h}\) 属于 \(\mathrm{E}_{\widetilde{m}_g}\) 且与 \(\mathrm{dom}(\widetilde{m}_{g'})\) 正交。这就意味着

\[
(\widetilde {h} \mid \widetilde {h} ^ {\prime}) _ {\widetilde {m} _ {g}} = \int_ {\mathrm{X}} \overline {{h}} h ^ {\prime} (1 + | g | ^ {2}) d \mu = 0
\]

它对每个这样的函数 \(h' \in \mathcal{L}_{\mathrm{K}}^{2}(\mathrm{X}, \mu)\) （其类 \(\widetilde{h}'\) 属于 \(\operatorname{dom}(\widetilde{m}_{g'})\) ）成立。

设 C 是 X 的紧子集， \(\varphi\) 是它的特征函数。设 \(n \in \mathbf{N}\) 。设 \(\varphi_n\) 是 \(\mu\)-可积集合 \(\mathrm{C}_n\) （由 C 中满足 \(x \in \mathrm{C}\) 的元素 \(|h(x)| \leqslant n\) 所成）的特征函数，并令 \(h_n' = \varphi_n h\) 。 \(h_n'\) 的类属于 \(\tilde{m}_{g'}\) 的定义域，因为 \(|g'h_n'| \leqslant n\varphi\) ；由此得出

\[
0 = \int_ {\mathrm{X}} \overline {{h}} h _ {n} ^ {\prime} (1 + | g | ^ {2}) d \mu = \int_ {\mathrm{X}} | h | ^ {2} \varphi_ {n} (1 + | g | ^ {2}) d \mu .
\]

这蕴含：对 \(\mu\)-几乎所有 \(x \in C_{n}\) ， h 为零；由于 n 任意，对 \(\mu\)-几乎所有 \(x \in C\) ， h 为零。最终推出 h \(\mu\) \(\mu\)-几乎处处为零，因为 C 任意且 h 是缓增的（INT, V, p. 9, § 1, n°3，推论）。

现在考虑断言 \(b\) )。由于 \(|g|^2\) 局部 \(\mu\)-可积，函数 \(fg\) 属于 \(\mathcal{L}_{\mathrm{K}}^{2}(\mathrm{X},\mu)\) （对 \(f\in\mathcal{K}(\mathrm{X};\mathrm{K})\) ），因此 \(\mathcal{K}(\mathrm{X};\mathrm{K})\) 包含在 \(\widetilde{m}_{g}\) 的定义域中。

设 \(h \in \mathcal{L}_{\mathrm{K}}^{2}(\mathrm{X}, \mu)\) ，其类 \(\widetilde{h}\) 属于 \(\mathrm{E}_{\widetilde{m}_{g}}\) 且与 G 正交。这时我们有

\[
0 = (\widetilde {h} \mid \widetilde {h} ^ {\prime}) _ {\widetilde {m} _ {g}} = \int_ {\mathrm{X}} h \overline {{h}} ^ {\prime} (1 + | g | ^ {2}) d \mu
\]

对所有 \(h' \in F\) ，其类为 \(h'\) 。鉴于关于 F 的假设，这意味着测度 \(h(1 + |g|^{2}) \cdot \mu\) 为零，因而 h \(\mu\)-几乎处处为零，因为 h 是缓增的。

设 \(p\) 是实数 \(\geqslant 1\) 。设 \(h\) 是 \(\mathcal{L}_{\mathrm{K}}^{\infty}(\mathrm{X},\mu)\) 中的元素。以 \(\widetilde{m}_h\) 为乘子的乘法算子 \(h\) 是 \(\mathrm{L}_{\mathrm{K}}^{p}(\mathrm{X},\mu)\) 的一个自同态。假设由 \(x\in X\) 中满足 \(h(x)=0\) 的 x 所成的集 Y 是局部 \(\mu\)-可忽略的。这时自同态 \(\widetilde{m}_h\) 是单射的（IV, p. 186 的引理 7）。用 \(h^{-1}\) 表示 X 上在 Y 上等于 0 、在 \(x \mapsto 1/h(x)\) 上等于 \(\mathrm{X} - \mathrm{Y}\) 的函数。逆部分算子 \(\widetilde{m}_{h}^{-1}\) 就是以 \(h^{-1}\) 为乘子的、从 \(\mathrm{L}_{\mathrm{K}}^{p}(\mathrm{X},\mu)\) 到 \(\mathrm{L}_{\mathrm{K}}^{p}(\mathrm{X},\mu)\) 中的乘法算子，即 \(\widetilde{m}_{h}^{-1} = \widetilde{m}_{h^{-1}}\) 。事实上， \(\widetilde{m}_{h}\) 的像是由形如 \(g \in \mathcal{L}_{\mathrm{K}}^{p}(\mathrm{X},\mu)\) （其中 \(g = hf\) ）的函数 \(f \in \mathcal{L}_{\mathrm{K}}^{p}(\mathrm{X},\mu)\) 的类所成的空间。这一条件等价于：对所有 \(g(x)/h(x) = f(x)\) 有 \(x \in \mathrm{X} - \mathrm{Y}\) ，且若 \(g(x) = 0\) 则 \(x \in \mathrm{Y}\) 。这蕴含： \(\widetilde{m}_{h}^{-1}\) 在 \(\mathrm{L}_{\mathrm{K}}^{p}(\mathrm{X},\mu)\) 中的定义域就是 \(\widetilde{m}_{h^{-1}}\) 的定义域，且等式 \(\widetilde{m}_{h}^{-1} = \widetilde{m}_{h^{-1}}\) 成立。

在下文中，我们有时把部分乘法算子简记为 \(m_{h}\) ，它以 h 为乘子，是从 \(\mathrm{L}_{\mathrm{K}}^{p_{1}}(\mathrm{X},\mu)\) 到 \(\mathrm{L}_{\mathrm{K}}^{p_{2}}(\mathrm{X},\mu)\) 中的。

\begin{enumerate}
\def\labelenumi{\arabic{enumi})}
\setcounter{enumi}{1}
\tightlist
\item
  设 E 是 K 上的希尔伯特空间， \(\mathrm{B} = (e_i)_{i \in \mathrm{I}}\) 是 E 的一个标准正交基。设 \((\lambda_i)_{i \in \mathrm{I}}\) 是 K 中元素的一个族。设 D 是 E 中由元素 \(x \in E\) （使得族 \((\lambda_i \langle e_i | x \rangle)_{i \in \mathrm{I}}\) 在 K 中平方可和）所成的向量子空间。空间 D 在 E 中稠密，因为它包含向量 \(e_i\) （对所有 \(i \in I\) ）。以 D 为定义域的部分算子 u 由下式给出
\end{enumerate}

\[
x \mapsto \sum_ {i \in I} \lambda_ {i} \langle e _ {i} | x \rangle e _ {i}
\]

称为基 B 中的部分对角算子，而 \((\lambda_{i})_{i\in I}\) 称为 u 的特征值族。

算子 \(u\) 是闭的。事实上，设 \((x_{n}, u(x_{n}))_{n \in \mathbf{N}}\) 是 \(u\) 的图像中在 \(\mathrm{E} \times \mathrm{E}\) 内收敛的一个元素序列，并设 \((x, y)\) 为其极限。这时我们有 \(\langle e_{i} | x_{n} \rangle \to \langle e_{i} | x \rangle\) （对所有 \(i \in \mathrm{I}\) ），并且

\[
\langle e _ {i} | u (x _ {n}) \rangle = \lambda_ {i} \langle e _ {i} | x _ {n} \rangle \rightarrow \langle e _ {i} | y \rangle
\]

对所有 \(i \in I\) 。于是， \(\lambda_{i} \langle e_{i} \mid x \rangle = \langle e_{i} \mid y \rangle\) 对所有 \(i \in I\) 成立；这就证明了 \(x \in D\) 且 \(u(x) = y\) ，即 \(u\) 是闭的。

这个例子实际上是前一例子应用于拓扑空间 X = I （赋以离散拓扑与计数测度 \(\mu\) ）的特殊情形，因为 E 通过映射 \(\ell^{2}(\mathrm{I}) = \mathrm{L}^{2}(\mathrm{I}, \mu)\) 等同于空间 \(x \mapsto (\langle e_{i} \mid x \rangle)_{i \in \mathrm{I}}\) （EVT, V, p. 23，推论 2），而 u 这时等同于乘法算子 \(m_{\lambda}\) ，其中 \(\lambda\) 是函数 \(i \mapsto \lambda_{i}\) 。

\begin{enumerate}
\def\labelenumi{\arabic{enumi})}
\setcounter{enumi}{2}
\tightlist
\item
  集合 \(N_{R} = N \cup \{\infty, \omega\}\) 赋有 VAR, R2, p. 10 中所述的全序，使得 \(n < \infty < \omega\) 对所有 \(n \in N\) 成立。设 \(r \in N_{R}\) 。设 \(n \in N\) ，并设 U 是 \(R^{n}\) 的一个开子集。设 \(k \in N\) 满足 \(k \leqslant r\) 。设 \((n_{\alpha})_{|\alpha| \leqslant k}\) 是 \(\mathcal{C}^{r}(U)\) 中元素的一个族，其中所考虑的重指标属于 \(N^{n}\) 。族 \((n_{\alpha})\) 在 U 上定义一个阶 \(\leqslant k\) 的纯量微分算子 D （参见 VAR, R2, 14.1.6, 14.1.4）。
\end{enumerate}

对每个满足 \(k \leqslant m \leqslant r\) 的整数 m ，微分算子 D 定义一个从 \(\mathrm{C}^{m}(\mathrm{U})\) 到 \(\mathrm{C}^{m-k}(\mathrm{U})\) 中的线性映射，它把 \(f \in \mathrm{C}^{m}(\mathrm{U})\) 映为

\[
\mathrm{D} (f) = \sum_ {| \alpha | \leqslant k} n _ {\alpha} \partial^ {\alpha} (f).
\]

同一公式定义了一个从 \(\mathcal{D}(\mathrm{U})\) 到 \(\mathcal{D}'(\mathrm{U})\) 中的连续线性映射（IV, p. 208 的定义 2）。

定义 5. --- 设 E 是 \(\mathcal{D}'(\mathrm{U})\) 的一个包含 \(\mathcal{D}(\mathrm{U})\) 的向量子空间。 E 上与 D 相伴的微分算子，指的是 E 上这样一个部分算子：它是以下述部分算子的扩张，该算子以 \(\mathcal{D}(\mathrm{U})\) 为定义域、由 \(\varphi \mapsto \mathrm{D}(\varphi)\) 所定义。

例如，设系数 \(n_{\alpha}\) 是 U 上的有界函数。设 \(\mu\) 是 U 上的勒贝格测度。若 p 是 \([1,+\infty[\) 中的一个元素，则可在 \(\mathrm{L}^{p}(\mathrm{U})\) 上定义一个与 D 相伴的微分算子，其定义域是索伯列夫空间 \(\mathrm{W}^{k,p}(\mathrm{U})\) （IV, p. 221 的第 14 节），因为这时我们有 \(n_{\alpha}\partial^{\alpha}(f)\in\mathrm{L}^{p}(\mathrm{U})\) （对所有 \(f\in\mathrm{W}^{k,p}(\mathrm{U})\) 与每个 \(|\alpha|\leqslant k\) ）。

